
\subsection{Stockage des données SAR embarquées}
La phase de stockage constitue un maillon critique de la chaîne d'acquisition SAR. Le radar génère un flux cohérent en phase, continu et de forte volumétrie, qu’il faut enregistrer sans perte ni désynchronisation. Le système de stockage doit donc combiner un débit soutenu, une grande fiabilité et une faible consommation énergétique.

\subsubsection{Rôle et contraintes générales}

Le stockage SAR embarqué a pour fonction principale de conserver l’ensemble des signaux I/Q bruts nécessaires à la formation d’image après la mission. Contrairement à d’autres capteurs (optiques ou lidar), les données SAR ne peuvent être ni sous-échantillonnées ni compressées fortement sans altérer la cohérence de phase. Les contraintes majeures sont les suivantes :
\begin{itemize}
  \item \textbf{Débit soutenu} : souvent supérieur à 10--100~MB/s selon la bande et le nombre de canaux.
  \item \textbf{Enregistrement ininterrompu} : aucune perte d’échantillons (``no~drop'') n’est tolérée.
  \item \textbf{Horodatage précis} : chaque bloc de données doit être lié à une référence temporelle (1~PPS, GNSS).
  \item \textbf{Intégrité} : absence d’erreurs d’écriture et cohérence du contenu dans le temps.
\end{itemize}

\subsubsection{Architecture typique}

\begin{center}
\texttt{[Antenne RX]} $\rightarrow$ \texttt{[Numériseur / FPGA]} $\rightarrow$ \texttt{[Buffer RAM]} $\rightarrow$ \texttt{[SSD / NVMe]} $\rightarrow$ \texttt{[Transfert post-vol]}
\end{center}

Le \textit{buffer} mémoire (DDR ou LPDDR) absorbe les variations de débit et protège contre les pertes de données, tandis que le stockage permanent (SSD ou NVMe) assure la persistance à long terme. Les fichiers sont généralement découpés par segments temporels (1--5~min) afin de limiter la corruption potentielle en cas d’arrêt inopiné.

\subsubsection{Technologies de stockage}

Les technologies les plus courantes pour les systèmes SAR embarqués sont résumées dans le tableau~\ref{tab:stockage_sar}.

\begin{table}[h!]
\centering
\caption{Comparatif des technologies de stockage pour applications SAR embarquées.}
\label{tab:stockage_sar}
\renewcommand{\arraystretch}{1.1}
\resizebox{\textwidth}{!}{%
\begin{tabular}{@{}lcccc@{}}
\toprule
\textbf{Technologie} & \textbf{Interface} & \textbf{Débit typique} & \textbf{Avantages} & \textbf{Inconvénients} \\
\midrule
SSD SATA III & 6~Gb/s & 300--500~MB/s & Fiable, économique & Débit limité, latence plus élevée \\
SSD NVMe PCIe 3/4 & 32~Gb/s & 1000--7000~MB/s & Très haut débit, faible latence & Échauffement, consommation accrue \\
eMMC / SDXC & -- & 30--100~MB/s & Compact, léger & Faible robustesse, pas pour flux continus \\
RAM circulaire & DDR4/5 & $>$5~GB/s & Très rapide, faible latence & Capacité limitée (tampon uniquement) \\
RAID embarqué & SATA/NVMe & variable & Redondance, fiabilité accrue & Masse et complexité supplémentaires \\
\bottomrule
\end{tabular}%
}
\end{table}

Les systèmes SAR légers (UAV~<~10~kg) privilégient aujourd’hui les SSD~NVMe PCIe~Gen~3/4, qui offrent un compromis entre débit (jusqu’à plusieurs~Go/s) et masse réduite. Le stockage sur mémoire~eMMC reste réservé aux systèmes de télémétrie ou à l’enregistrement d’état, mais non au signal radar brut.

\subsubsection{Formats de fichiers et structuration}

Les données sont classiquement enregistrées sous forme binaire, composées de couples d’échantillons complexes~(I/Q). Chaque fichier contient un en-tête décrivant les paramètres d’acquisition : fréquence~$f_0$, bande~$B$, fréquence d’échantillonnage~$f_s$, horodatage, position~GNSS, polarisation, et température électronique.

Depuis 2017, le format \textbf{SigMF} (\textit{Signal Metadata Format}) est devenu un standard de fait pour l’archivage de signaux RF~\cite{Balister2017_SigMF}. Il associe un fichier de données brutes (\texttt{.sigmf-data}) et un fichier de métadonnées JSON (\texttt{.sigmf-meta}), assurant traçabilité et interopérabilité avec les environnements GNU~Radio, MATLAB ou Python.

\subsubsection{Réduction et compression de données}

Différentes approches sont étudiées pour réduire la volumétrie des données SAR tout en préservant la cohérence de phase :
\begin{itemize}
  \item \textbf{Compression sans perte} : codage adaptatif ou entropique avec des gains de 1.3 à 2× selon la structure du signal~\cite{Pieterse2019_SARcompressionMetrics}.
  \item \textbf{Compression avec perte contrôlée} : méthodes basées sur la transformée en cosinus ou les représentations par blocs compressifs~\cite{Choi2021_DroneSARcompression, Kim2025_SARcompressionPointClouds}.
  \item \textbf{Décimation contrôlée} après déchirpage lorsque la bande utile est inférieure à $f_s$.
  \item \textbf{Prétraitement embarqué} (aperçu \textit{range--Doppler}) pour visualisation rapide sans transfert complet.
\end{itemize}

Ces méthodes sont employées avec précaution dans les systèmes UAV, la préservation de la phase complexe restant essentielle à la focalisation SAR~\cite{Svedin2021_UAVSAR}.

\subsubsection{Aspects énergétiques et fiabilité}

Le stockage haute vitesse induit une consommation non négligeable : un SSD~NVMe en écriture soutenue consomme environ~1.5--2~W. La sous-chaîne \textit{acquisition + stockage} représente typiquement~10--20\,\% de la puissance totale d’une charge utile SAR~\cite{Sun2020_PassiveUAVSAR}. Les bonnes pratiques incluent :
\begin{itemize}
  \item l’usage d’un système de fichiers journalisé (ext4, XFS) en mode \texttt{noatime};
  \item un \textit{flush} périodique des buffers et un condensateur de sauvegarde pour terminer l’écriture en cas de coupure;
  \item le calcul d’un hash ou CRC par segment et la génération d’un fichier manifeste global.
\end{itemize}

\subsubsection{Tendances récentes}

Les recherches actuelles explorent l’utilisation de mémoires non volatiles résistives (ReRAM, MRAM) pour des stockages plus robustes en environnement vibrant, ainsi que des architectures \textit{edge computing} combinant compression adaptative et détection de cibles à bord. L’objectif est de réduire la dépendance à la mémoire de masse tout en préservant la cohérence nécessaire à la synthèse d’ouverture.

\subsubsection{Résolution de quantification nécessaire pour la précision de phase}

Dans un radar à synthèse d’ouverture, la fidélité de la \emph{phase complexe} conditionne la cohérence de la synthèse et, in fine, la résolution azimutale.
On modélise la quantification $I/Q$ par un bruit additif sur chaque composante :
\[
z = A\,e^{j\varphi} + n_I + j\,n_Q,\qquad n_I,n_Q \sim \mathcal{U}\!\left(-\tfrac{\Delta}{2},\tfrac{\Delta}{2}\right),
\]
où $\Delta$ est le pas de quantification et $\mathrm{Var}(n_I)=\mathrm{Var}(n_Q)=\sigma_q^2=\Delta^2/12$.
Pour des bruits faibles, l’écart type de phase introduit par la quantification vérifie
\[
\sigma_\varphi \approx \frac{\Delta}{\sqrt{12}\,A}
\quad\text{et}\quad
\sigma_\varphi \approx \frac{1}{\sqrt{2\,\mathrm{SNR}_{q,\mathrm{lin}}}},
\]
avec $\mathrm{SNR}_{q,\mathrm{lin}} \triangleq A^2/(2\sigma_q^2)$. Pour un CAN idéal à $N$ bits et un \emph{backoff} (marge anti-saturation) de $\mathrm{backoff}_{\mathrm{dB}}$,
\[
\mathrm{SNR}_{q,\mathrm{dB}} \simeq 6{,}02\,N + 1{,}76 - \mathrm{backoff}_{\mathrm{dB}}.
\]

\paragraph{Lien avec la contrainte \texorpdfstring{$\lambda/100$}{lambda/100}.}
À 24\,GHz, la longueur d’onde vaut $\lambda \approx 12{,}5$\,mm. Une erreur radiale $\Delta R=\lambda/100 \approx 0{,}125$\,mm provoque en monostatique une erreur de phase
\[
\Delta\varphi \approx \frac{4\pi\,\Delta R}{\lambda} \approx 0{,}126~\text{rad}~(7{,}2^\circ).
\]
Si l’on adopte une marge plus stricte en limitant l’écart type à $\sigma_\varphi \le 0{,}063$\,rad (3,6°), il suffit d’assurer $\mathrm{SNR}_{q,\mathrm{dB}}\gtrsim 21$\,dB pour que la \emph{seule} quantification ne limite pas la cohérence.

\paragraph{Exemple (I/Q sur 16 bits, backoff 12 dB).}
Avec $\mathrm{backoff}=12$\,dB, on obtient :
\begin{table}[H]
\centering
\caption{Erreur de phase RMS due à la quantification selon la résolution du CAN (24\,GHz, backoff 12\,dB).}
\label{tab:bits_phase_24ghz}
\renewcommand{\arraystretch}{1.1}
\begin{tabular}{@{}cccc@{}}
\toprule
$N$ (bits) & $\mathrm{SNR}_{q,\mathrm{dB}}$ & $\sigma_\varphi$ (rad) & $\sigma_\varphi$ (deg) \\
\midrule
8  & 37.9 & $9.0\times10^{-3}$  & $0.52^\circ$ \\
10 & 50.0 & $2.25\times10^{-3}$ & $0.13^\circ$ \\
12 & 62.0 & $5.62\times10^{-4}$ & $0.032^\circ$ \\
14 & 74.0 & $1.40\times10^{-4}$ & $0.008^\circ$ \\
\bottomrule
\end{tabular}
\end{table}

\noindent\textbf{Conclusion.} Même pour une exigence de l’ordre de $\lambda/100$, 8 à 10 bits \emph{effectifs} suffisent du point de vue de la quantification pure. En pratique, on retient des CAN 12–14 bits (ENOB $\gtrsim$ 10) pour garantir la dynamique de scène, la polarimétrie et la robustesse vis-à-vis de la gigue d’horloge, du bruit de phase LO et des erreurs de trajectoire.

\vspace{1em}

\subsubsection{Exemple représentatif : volume de données d’un vol à 24\,GHz}

Le flux de données généré par un radar SAR dépend directement des paramètres d’acquisition et du mode d’enregistrement.  
On considère ici un radar FMCW déchirpé, enregistrant des échantillons complexes $I/Q$ sur 16\,bits par composante (soit 4\,octets par échantillon et par canal).  
Les relations suivantes permettent d’estimer le débit et le volume total de données :
\[
d = \mathrm{PRF}\,T_{\mathrm{chirp}}, \qquad
R_{\mathrm{inst}} = 4\,f_s\,N_c, \qquad
R_{\mathrm{moy}} = d\,R_{\mathrm{inst}}, \qquad
V = R_{\mathrm{moy}}\,T_{\mathrm{vol}},
\]
relations standards pour les radars à synthèse d’ouverture~\cite{cumming2005,richards2010,Aguasca2013_ARBRES_UAV_SAR}.  
Elles s’interprètent de la manière suivante :
\begin{itemize}
  \item \textbf{$T_{\mathrm{chirp}}$}~: durée d’un signal élémentaire émis (\emph{chirp}) ; elle fixe la bande instantanée utile et donc la résolution en distance ($\delta_r = c / (2B)$).
  \item \textbf{$\mathrm{PRF}$}~: fréquence de répétition des impulsions (\emph{Pulse Repetition Frequency}) ; elle détermine la cadence d’acquisition azimutale et la portée non ambiguë.
  \item \textbf{$d = \mathrm{PRF}\,T_{\mathrm{chirp}}$}~: rapport cyclique (\emph{duty-cycle}) décrivant la fraction du temps d’émission par rapport à un cycle émission–écoute.
  \item \textbf{$f_s$}~: fréquence d’échantillonnage après déchirpage (bande de base) ; elle définit le nombre d’échantillons collectés par chirp.
  \item \textbf{$N_c$}~: nombre de canaux indépendants (souvent polarimétriques : HH, VV, etc.) ; chaque voie produit un flux $I/Q$ distinct.
  \item \textbf{$R_{\mathrm{inst}} = 4\,f_s\,N_c$}~: débit instantané en octets par seconde à la sortie des ADC (4 octets = 2 composantes $\times$ 2 octets par composante).
  \item \textbf{$R_{\mathrm{moy}} = d\,R_{\mathrm{inst}}$}~: débit moyen réellement enregistré, tenant compte du rapport cyclique $d$.
  \item \textbf{$T_{\mathrm{vol}}$}~: durée utile d’enregistrement pendant le vol (souvent quelques dizaines de minutes).
  \item \textbf{$V = R_{\mathrm{moy}}\,T_{\mathrm{vol}}$}~: volume total de données brutes enregistrées sur le support mémoire embarqué.
\end{itemize}

Enfin, la contrainte d’échantillonnage azimutal relie la PRF à la vitesse du porteur~\cite{cumming2005,richards2010} :
\[
\Delta x = \frac{v}{\mathrm{PRF}},
\]
où $v$ est la vitesse du drone et $\Delta x$ l’espacement spatial entre deux positions successives de l’antenne le long de la trajectoire.  
Pour éviter les ambiguïtés d’azimut, ce pas doit rester inférieur à $\lambda/2$.

\paragraph{Paramètres de référence (projet).}
Pour un système SAR compact en bande~K ($f_0=24$\,GHz, $\lambda\simeq12{,}5$\,mm) embarqué sur un drone léger :
\begin{itemize}
  \item vitesse moyenne du porteur : $v \approx 18$\,m/s (environ 65\,km/h) ;
  \item nombre de voies : $N_c=1$ ou $2$ selon la polarisation ;
  \item durée de vol utile : $T_{\mathrm{vol}}=20$ à $25$\,minutes.
\end{itemize}
Les trois profils ci-dessous représentent des scénarios réalistes :

\begin{table}[H]
\centering
\caption{Volume de données selon trois profils de mission (24\,GHz, $I/Q$ sur 16\,bits).}
\label{tab:volume_donnees_vol_24ghz}
\renewcommand{\arraystretch}{1.1}
\begin{tabular}{@{}lcccccccc@{}}
\toprule
\textbf{Profil} & $f_s$ & $T_{\mathrm{chirp}}$ & PRF & $d$ & $N_c$ & $R_{\mathrm{inst}}$ & $R_{\mathrm{moy}}$ & $V$ sur $T_{\mathrm{vol}}$ \\
 & (MS/s) & ($\mu$s) & (Hz) &  &  & (MB/s) & (MB/s) &  \\ \midrule
A — Léger   & 3  & 120 & 800  & 0.096 & 1 & 12.0 & 1.15 & \textbf{1.4\,GB} (20\,min) \\
B — Nominal & 6  & 200 & 1500 & 0.30  & 2 & 48.0 & 14.4 & \textbf{21.6\,GB} (25\,min) \\
C — Riche   & 10 & 250 & 2000 & 0.50  & 2 & 80.0 & 40.0 & \textbf{48\,GB}  (20\,min) \\
\bottomrule
\end{tabular}
\end{table}

\noindent\textit{Échantillonnage azimutal et cohérence.}  
Pour une vitesse $v=18$\,m/s :
\[
\Delta x = \frac{v}{\mathrm{PRF}} \Rightarrow
\begin{cases}
\text{A : } 2{,}25~\mathrm{cm},\\
\text{B : } 1{,}2~\mathrm{cm},\\
\text{C : } 9~\mathrm{mm},
\end{cases}
\]
ce qui reste compatible avec une antenne compacte à 24\,GHz ($\lambda\simeq12{,}5$\,mm).

\paragraph{Commentaires.}
\begin{itemize}
  \item Le \textbf{profil~B (nominal)} retient $f_s=6$\,MS/s, $N_c=2$ (dual-pol) et $d=0{,}30$. On obtient un débit moyen $R_{\mathrm{moy}}\approx14{,}4$\,MB/s, soit \textbf{21{,}6\,GB} pour 25\,min d’acquisition.
  \item Le \textbf{profil~C} (bande utile et PRF plus élevées) atteint $\sim48$\,GB en 20\,min. Un SSD~NVMe de 512\,GB couvre ainsi plusieurs missions complètes.
  \item Les valeurs indiquées concernent le \textbf{flux brut $I/Q$}. Un flux d’\emph{aperçu} (image range–Doppler en \texttt{float16}) ne représente que quelques centaines de~MB et permet le contrôle rapide en vol.
\end{itemize}

\vspace{0.5em}
En résumé, les équations précédentes relient directement les paramètres physiques d’acquisition ($f_s$, PRF, $T_{\mathrm{chirp}}$, $N_c$) aux besoins en débit et en capacité mémoire. Ces ordres de grandeur démontrent que, pour un radar FMCW embarqué en bande~K, le stockage du flux brut sur SSD~NVMe ($\geq$512\,GB) reste compatible avec une mission UAV de 20–30\,minutes tout en préservant la cohérence temporelle et la précision de phase~\cite{Aguasca2013_ARBRES_UAV_SAR,cumming2005,richards2010}.
En résumé, la résolution du CAN fixée par l’exigence de phase ($\lambda/100$) reste peu contraignante face aux autres sources d’erreurs (jitter d’horloge, bruit de phase LO, trajectoire). En revanche, $f_s$, PRF, $N_c$ et $d$ déterminent directement le débit moyen et, partant, la \textbf{capacité de stockage} et la \textbf{gestion de puissance} requises pour la mission.


\subsection*{Conclusion générale sur l'acquisition et le stockage SAR embarqué}

L’analyse des contraintes d’acquisition et de stockage met en évidence la cohérence globale de la chaîne SAR embarquée : chaque paramètre physique — fréquence d’échantillonnage, PRF, durée de chirp, nombre de canaux — influe directement sur le débit, la précision de phase et la capacité mémoire requise.

À 24\,GHz, la courte longueur d’onde impose une rigueur accrue sur la cohérence temporelle, la stabilité d’horloge et la qualité du front-end RF, mais la quantification numérique (12–14\,bits effectifs) ne constitue pas un facteur limitant. Les calculs montrent que la contrainte de phase de l’ordre de $\lambda/100$ est largement tenue, et que la dynamique disponible permet de préserver la qualité radiométrique et polarimétrique de l’imagerie.

Les estimations de débit et de volume de données confirment la faisabilité d’un enregistrement intégral du flux $I/Q$ sur mémoire SSD~NVMe pour des missions UAV de 20 à 30\,minutes. Un tel dispositif assure la cohérence nécessaire à la formation d’image tout en restant compatible avec les contraintes de masse et d’énergie d’un radar compact en bande~K.

En conclusion, la maîtrise conjointe de la quantification, de la cadence d’acquisition et de l’architecture de stockage constitue le socle matériel du système SAR envisagé. Ces fondations garantissent la qualité des données brutes et conditionnent les performances des traitements ultérieurs (focalisation, calibration, compression et géoréférencement).